% Experiment report. Fixed sections: see experiments/TEMPLATE/report.tex and experiments/README.md.
% Build: npm run experiment:report -- 26-rl-spend4-vi-r15
\documentclass{../polyexp}
\begin{document}
% generated by scripts/experiment-report.ts from experiment.json and models/<run>/ - do not edit
\providecommand{\ExpHeader}{}\renewcommand{\ExpHeader}{%
{\sffamily\footnotesize\color{expgrey}EXPERIMENT 26 \textperiodcentered\ METHOD \hfill 2026-11-27}\par
{\LARGE\bfseries\sffamily Spending actions with the unchanged bot as an action, value iteration on r15}\hfill\verdictbox{expred}{NOT KEPT}\par\smallskip
{\itshape If the unchanged bot is one of the actions (bot, economy, army, research), does value iteration on r15 find turns where a spending boost pays, or does it learn to leave the bot alone?}\par\medskip
\begin{tabular}{@{}p{22mm}p{\dimexpr\linewidth-26mm}@{}}
\fact{Change}{Experiment 25 (state with the turn) with the action set spend4 = bot, economy, army, research instead of spend3; the bot is also the default for states without an entry.}
\fact{Knobs}{\knob{minCount} $\in$ \{1, 5, 20, 50, 200\}, best 200}
\fact{Base}{\model{greedy-full-probe-perf-r15} $\rightarrow$ result \model{rl-gate vs greedy-full-probe-perf-r15 FAIL (-210, 43.7\% better); vs experiment 25 policy +30, not significant; random placement of the same boosts costs only about -45}}
\fact{Method}{300k games with a random spend4 action per turn; value iteration with the turn and min-count 1, 5, 20, 50, 200; min-count chosen on 2,000 validation seeds; the chosen policy on the 10,000 evaluation seeds against r15 and against experiment 25's policy; metric: final raw score, \#171 rule}
\fact{Setups}{Perfection, 16$\times$16 continents, 4 players, Crazy, seat 0 Xin-Xi (human seat), learner base greedy-full-probe-perf-r15; opponents greedy-full-probe-perf-r15, greedy-full-probe-perf, greedy-full-probe-perf (anchor), seats shuffled per seed (rl-gate defaults from 27 Nov)}
\fact{Record}{issue \issue{166}; status: done}
\end{tabular}\par\medskip
\begin{center}\begin{tikzpicture}[node distance=5mm]
    \node[flowbox] (n0) {{\sffamily\bfseries\scriptsize\color{expgrey}BASE}\\[1pt]\model{full-probe-perf-r15}};
    \node[flowbox,right=5mm of n0] (n1) {{\sffamily\bfseries\scriptsize\color{expgrey}CHANGE}\\[1pt]\knob{minCount}\\{\scriptsize 1, 5, 20, 50, 200}};
    \node[flowbox,right=5mm of n1] (n2) {{\sffamily\bfseries\scriptsize\color{expgrey}MEASURED}\\[1pt]Paired, 10,000 evaluation seeds\\{\scriptsize 1 setup, metric: final raw score, \#171 rule}};
    \node[flowbox,right=5mm of n2] (n3) {{\sffamily\bfseries\scriptsize\color{expgrey}RESULT}\\[1pt]Mean final score, 10,000 evaluation seeds: r15 vs spend4 policy (m = 200)\\{\scriptsize 7199.0 $\rightarrow$ 6989.0}\\\textbf{$-$210.0}};
    \node[flowbox,right=5mm of n3] (n4) {{\sffamily\bfseries\scriptsize\color{expgrey}VERDICT}\\[1pt]\textbf{\color{expred}NOT KEPT}\\{\scriptsize new champion:}\\\model{rl-gate vs greedy-full-probe-perf-r15 FAIL (-210, 43.7\% better); vs experiment 25 policy +30, not significant; random placement of the same boosts costs only about -45}};
    \draw[flowarrow] (n0) -- (n1); \draw[flowarrow] (n1) -- (n2); \draw[flowarrow] (n2) -- (n3); \draw[flowarrow] (n3) -- (n4);
\end{tikzpicture}\end{center}
\begin{center}\sffamily\small\begin{tabular}{@{}lrrr@{}}\toprule
Measure & Base & New & Change\\\midrule
Mean final score, 10,000 evaluation seeds: r15 vs spend4 policy (m = 200) & 7199.0 & 6989.0 & $-$210.0\\
Same seeds: experiment 25 policy vs spend4 policy & 6959.0 & 6989.0 & +30.0\\
\bottomrule\end{tabular}\end{center}
}
\providecommand{\ExpResults}{}\renewcommand{\ExpResults}{%
{\small\color{expgrey}No runs recorded.}
}

\ExpHeader

\section{Question and hypothesis}
\textbf{The idea (David, 27 Nov 2026, issue \#166).} Experiment 25 again, with the unchanged bot as a fourth action:
\texttt{spend4} = bot, economy, army, research. In 25 the action set had no ``do nothing'' action, so every policy
decision was a boost; the noisy maximum over three $Q$ estimates chose economy on 16.8\% of learner turns, and the
policy lost 240 points against r15. With the bot as an action, a boost is played only where its estimate beats the
unchanged bot in that state. This asks whether any turn-dependent boost pays on r15 at all.

\textbf{Prediction} (before any run of this experiment):
\begin{enumerate}
  \item \emph{The policy leaves the bot alone more often.} At the chosen min-count, the bot plays unchanged (as the
        chosen action or by fallback) on at least 40\% of learner turns, and economy on fewer than 10\%
        (25: bot 30.5\%, economy 16.8\%).
  \item \emph{Still no gain over r15.} The mean difference against r15 on the evaluation seeds is between $-150$ and
        $+30$ points, and the policy does \emph{not} pass the \#171 rule against r15.
  \item \emph{Better than experiment 25.} Paired on the evaluation seeds, the policy beats 25's policy (with turn,
        $m = 50$) by at least $+80$ points, and the sign test rejects ``no better'' at $\alpha = 0.01$.
  \item \emph{The larger the min-count, the better.} On validation, $m = 200$ is best or within noise of the best
        (fewer qualifying actions, more turns played by the bot).
\end{enumerate}
\textbf{Against the prediction:} a PASS against r15 would mean that value iteration finds timing decisions that the
research-corrected bot misses. A policy that is no better than 25's would mean the loss in 25 does not come from
the missing bot action.

\section{Mechanism}
\textbf{Actions} (\texttt{spend4} in \texttt{rl\_value::strategic}, new): 0 bot (the base weights unchanged),
1 economy (build), 2 army (train, attack), 3 research; each boost multiplies the base bot's move-type factors of its
types by 2 for one learner turn, as in 23 and 25. Action 0 is also the default: a state without a policy entry
plays the bot unchanged. A test checks that action 0 is exactly r15 and actions 1--3 are spend3's.

\textbf{State}: (turn, cities, income, units, techs), as the chosen version of 25 (without the turn lost more there).

\textbf{Value iteration} as in experiment 25 ($Q(s,a) = $ mean of $r + V(s')$, $V(s) = \max_{a: n(s,a) \ge m}
Q(s,a)$, fallback to the mean over all transitions, $\gamma = 1$), one backward sweep from the last turn.

\section{Setup}
As experiment 25: learner base \texttt{greedy-full-probe-perf-r15}; opponents r15, probe-perf and the anchor
probe-perf (\texttt{rl-gate} defaults after PR \#204), shuffled per seed. Perfection, 16$\times$16 Continents,
4 players, Crazy, learner Xin-xi in the human seat. Seeds: training $1\dots300{,}000$ (new games, random spend4
action per turn, so the bot is played on a quarter of the turns), action check $600{,}000{,}000\dots$ (2{,}000),
validation $700{,}000{,}000\dots$ (2{,}000), evaluation $10^9\dots$ (10{,}000). Min-count $1, 5, 20, 50, 200$, chosen
on validation; the chosen policy on the evaluation seeds with \texttt{rl-gate} against r15 and against 25's policy
(\texttt{--base rl-policy:experiments/25-rl-spend-vi-r15/policy-turn-m50.tsv}).

Commands (from \texttt{engine/}):
\begin{verbatim}
B=./target/release/rl-value; R=target/rl-value; E=../experiments/26-rl-spend4-vi-r15
$B check --actions spend4 --games 2000 --seed-start 600000000 --out $E/check.json
$B collect --actions spend4 --driver random --games 300000 --seed-start 1 --out $R/s4r15-train-random.tsv
$B fit-vi --train $R/s4r15-train-random.tsv --features turn,cities,income,units,techs \
  --min-count 1,5,20,50,200 --out $R/vi-s4r15-turn
$B eval --actions spend4 --policy $R/vi-s4r15-turn/policy-mM.tsv --n 2000 --seed-start 700000000 --out $E/val-mM.json
rl-gate --candidate rl-policy:experiments/26-rl-spend4-vi-r15/policy-mM.tsv --base greedy-full-probe-perf-r15
rl-gate --candidate rl-policy:experiments/26-rl-spend4-vi-r15/policy-mM.tsv \
  --base rl-policy:experiments/25-rl-spend-vi-r15/policy-turn-m50.tsv
\end{verbatim}

\section{Results}
\ExpResults

All numbers are copied from \texttt{check.json}, \texttt{vi-summary-turn.json}, \texttt{val-*.json},
\texttt{eval-m200.json} and the two \texttt{rl-gate} JSONs in this folder (27 Nov 2026; per-seed rows dropped, the
runs are deterministic). Data: 300{,}000 games with a random spend4 action per turn on r15, 8{,}463{,}276 rows,
1{,}526\,s on 10 threads; 260{,}583 distinct states.

\paragraph{Action check} (2{,}000 seeds, each action every turn): the same games as in 25 (same seeds, base and
opponents): bot 7{,}148, economy 6{,}687 ($-461$), army 7{,}157 ($+9$), research 7{,}112 ($-36$).

\paragraph{Min-count on validation} (2{,}000 seeds against r15, base mean 7{,}282). ``Bot'' counts turns where the
policy chose the bot and turns without an entry; $V(\text{start})$ is value iteration's own estimate of the expected
final score.

\begin{center}\small
\begin{tabular}{rrrrrrrrr}
\toprule
$m$ & entries cover & $V(\text{start})$ & mean $d$ & better / worse / ties & bot & economy & army & research \\
\midrule
1   & 100.0\% & 9{,}705 & $-313$ & 818 / 988 / 194 & 26.0\% & 25.8\% & 25.5\% & 22.8\% \\
5   & 93.0\%  & 8{,}230 & $-213$ & 853 / 944 / 203 & 31.5\% & 21.2\% & 25.7\% & 21.6\% \\
20  & 78.9\%  & 7{,}542 & $-210$ & 838 / 918 / 244 & 40.4\% & 17.8\% & 23.9\% & 17.9\% \\
50  & 65.1\%  & 7{,}171 & $-225$ & 782 / 872 / 346 & 50.1\% & 15.0\% & 20.9\% & 14.1\% \\
200 & 43.7\%  & 6{,}821 & $-187$ & 685 / 764 / 551 & 65.4\% & 11.2\% & 14.0\% & 9.3\% \\
\bottomrule
\end{tabular}
\end{center}

Every min-count loses against r15; $m = 200$ is chosen (best mean, but within noise of 5 and 20).

\paragraph{Evaluation seeds} (10{,}000, \texttt{rl-gate}):
\begin{itemize}
  \item \emph{Against r15: FAIL.} Policy 6{,}989 vs r15 7{,}199, mean $d = -210$ (upper bound $-158$); better on
        3{,}141, worse on 4{,}044, tied on 2{,}815 seeds (43.7\% of differing seeds better); p10 1{,}435 vs
        1{,}480. 15 opponents (information only): $-276$, 40.9\% better.
  \item \emph{Against experiment 25's policy (with turn, $m = 50$): no significant difference.} 6{,}989 vs 6{,}959,
        mean $d = +30$ (99\% bounds $-4$ to $+64$); better 3{,}565, worse 3{,}481 (50.6\%, $p = 0.16$).
\end{itemize}

\section{Behaviour}
On the evaluation seeds the policy plays the bot on 65.1\% of learner turns (chosen or no entry; 25: 30.5\%),
economy 11.4\% (25: 16.8\%), army 14.1\%, research 9.4\%.

\paragraph{Where the boosts fall.} Weighted by visits, the policy boosts on 100\% of turns 0--3 and on 59--77\% of
turns 4--11 in states with an entry. The most visited states have almost equal $Q$ values:

\begin{center}\small
\begin{tabular}{lrrrrrl}
\toprule
state (turn, cities, income, units, techs) & visits & $Q$ bot & economy & army & research & chosen \\
\midrule
0, 1, 2, 1, 1 (start) & 300{,}000 & 6{,}728 & 6{,}728 & 6{,}727 & 6{,}821 & research \\
2, 1, 2, 2, 2         & 221{,}084 & 6{,}336 & 6{,}322 & 6{,}322 & 6{,}336 & research \\
1, 1, 2, 2, 2         & 176{,}708 & 6{,}727 & 6{,}728 & 6{,}725 & 6{,}724 & economy \\
3, 1, 2, 2, 2         & 137{,}214 & 5{,}826 & 5{,}835 & 5{,}818 & 5{,}826 & economy \\
\bottomrule
\end{tabular}
\end{center}

\paragraph{Shuffle test} (added after the gate, \texttt{shuffle.py}): the policy's actions permuted at random over its
1{,}025 entries (same states, about the same shares, no link between state and action), on the validation seeds:

\begin{center}\small
\begin{tabular}{lrrrr}
\toprule
policy & mean $d$ vs r15 & better / worse / ties & mean $d$ on differing seeds & bot share \\
\midrule
value iteration, $m = 200$ & $-187$ & 685 / 764 / 551   & $-258$ & 65.4\% \\
shuffled, seed 1           & $-51$  & 307 / 339 / 1{,}354 & $-158$ & 72.4\% \\
shuffled, seed 2           & $-39$  & 335 / 313 / 1{,}352 & $-119$ & 68.1\% \\
\bottomrule
\end{tabular}
\end{center}

\section{Interpretation}
\textbf{Prediction 1 (bot on at least 40\%, economy below 10\%): half right.} The bot plays 65.1\% of turns, but
economy 11.4\%.
\textbf{Prediction 2 ($-150$ to $+30$, no PASS): no PASS right, the range wrong} ($-210$).
\textbf{Prediction 3 (at least $+80$ better than 25, significant): wrong} ($+30$, $p = 0.16$).
\textbf{Prediction 4 (larger min-count better): right}, $m = 200$ is best on validation, though within noise of
5 and 20.

\textbf{The bot as an action did not help, because the policy rarely leaves it where it matters.} The same number of
boosts placed at random costs about 45 points; placed by value iteration, 187 (validation) and 210 (evaluation). Two
effects:
\begin{enumerate}
  \item \emph{The opening.} In the common early states the four $Q$ values are within a few points of each other.
        Three boosts compete against one bot, so the maximum picks a boost about three times in four when all are
        equal, and on turns 0--3 it picks one every time. Nearly every game leaves the bot's path in its first turns
        (551 ties against 1{,}350 for the shuffled policies), and r15 is near a local optimum of its weights, so
        small deviations cost on average.
  \item \emph{Worse than chance where it deviates} ($-258$ per differing game against $-119$ to $-158$ shuffled).
        $Q(s,a) = r + V(s')$ with $V(s')$ a maximum over noisy estimates: actions whose successors were estimated
        high, often rare and noisy states, are preferred, and those states are often rare because they are bad.
        Value iteration's own estimates show the optimism: at $m = 1$ it expects 9{,}705 points from the start state,
        the policy scores 6{,}968. At the start state the gap research $-$ bot ($+93$) is larger than its sampling
        error (about 15 points at 75{,}000 games per action), so it comes from the successors' maxima, not from the
        start state's own noise.
\end{enumerate}
Together with 25, this says the limit is the decision rule on near-equal noisy estimates, not the missing action.

\section{Threats to validity}
\begin{itemize}
  \item \emph{The shuffle test is two shuffles on 2{,}000 seeds}, added after the gate; it shows the size of the
        effect, not a precise number.
  \item \emph{Data from random actions.} $Q$ is estimated from games where every turn's action is random; the
        policy's own state distribution differs, and with an aliased state (stars, free unit slots and threats are
        not in it) the transitions under the policy are not those of the data.
  \item \emph{One boost size} (×2) and one state version; a smaller boost would make deviations cheaper.
\end{itemize}

\section{Verdict and next step}
\textbf{Not kept}: FAIL against r15 ($-210$), and not better than 25. r15 stays the base; David's RL v2 candidate
is still experiment 23's policy (measured on the old base bot).

Next step (one idea, same data, no new games): a decision rule that leaves the bot alone unless a boost is clearly
better: boost only if $Q(s, \text{boost}) - Q(s, \text{bot})$ exceeds about 2--3 standard errors, or double
estimation (choose the action on one half of the games, estimate its value on the other, as in double
$Q$-learning). After that: a coarse, decision-relevant state (free unit slots, stars, threatened city, unclaimed
villages, score gap).

\section{Notes for the paper}
Draft claim: on a well-tuned base bot, tabular value iteration over one-turn weight boosts loses 210 points, and
a random placement of the same boosts loses only about 45: the maximum over near-equal noisy $Q$ values makes the
policy deviate in every game's opening and favours successors whose values are overestimated. Carried by the
shuffle-test table.
\todo{jakob}{Is this worth a place in the paper? Where?}

\end{document}
